% !TEX program = xelatex
% 컴파일: xelatex squeezing_report_v4.tex  (목차/참조 때문에 2회 실행 권장)
% 한글 출력: kotex + XeLaTeX. Windows 기본 'Malgun Gothic' 사용.
% 그림:
%   squeezing_frontier_v4.png            -- v4 finite-loss ultra frontier (교정된 물리)
%   squeezing_frontier_od_v4.png         -- v4 in-cell OD 진단
%   squeezing_frontier_ideal_qe1_v4.png  -- v4 ideal QE=1, loss=0 wide dense scan (교정된 물리)
\documentclass[11pt,a4paper]{article}

\usepackage{kotex}
\setmainhangulfont{Malgun Gothic}
\setsanshangulfont{Malgun Gothic}

\usepackage{amsmath,amssymb}
\usepackage{graphicx}
\usepackage{booktabs}
\usepackage{array}
\usepackage{geometry}
\geometry{margin=2.4cm}
\usepackage{caption}
\captionsetup{font=small,labelfont=bf}
\usepackage{enumitem}
\usepackage{xcolor}
\usepackage[hidelinks]{hyperref}
\graphicspath{{./}}
\sloppy

\newcommand{\dB}{\,\mathrm{dB}}
\newcommand{\GHz}{\,\mathrm{GHz}}
\newcommand{\MHz}{\,\mathrm{MHz}}
\newcommand{\degC}{\,^{\circ}\mathrm{C}}

\title{\textbf{$^{85}$Rb D1 이중-$\Lambda$ 사파장혼합(FWM) 트윈빔\\
세기차 스퀴징 최적점: 이상 노이즈 공식 교정과\\
충돌 결어긋남 도입 이후의 재평가}\\[4pt]
\large 보고서 v4}
\author{GABES (Generic Atomic Bloch Equation Solver) 분석 보고서}
\date{2026년 7월 1일}

\begin{document}
\maketitle

\begin{abstract}
\noindent
v3 이후 GABES FWM 엔진에 두 가지 물리 교정이 들어갔다. 첫째(A), 이상(무손실)
트윈빔 세기차 노이즈 공식을 $(G_s+G_c-2\sqrt{G_sG_c-1})/(G_s+G_c)$에서 보존량
논증에 따른 정확식 $S_{\mathrm{ideal}}=(G_s-G_c)^2/(G_s+G_c)$(무손실 한계에서
$1/(2G-1)$로 환원)로 교정했다. 둘째(E), 바닥상태 및 광학 결맞음에 밀도·온도
의존 충돌 결어긋남(collisional decoherence)을 도입해, 고온에서 squeezing이
무한정 개선되지 않고 실제로 악화되는 물리를 반영했다.

본 v4 보고서는 이 두 교정이 v3의 세 가지 결과 — (1) 실험 레퍼런스 재현,
(2) 현실적 손실 포함 최효율 전선, (3) 이상 검출 한계 스캔 — 에 미친 영향을
재평가한다. \textbf{핵심 발견은 v3의 이상 스캔 결과($\xi=-81.9469\dB$)가 물리적
실체가 아니라 교정 전 공식의 버그였다는 것이다.} 같은 $(G_s,G_c)\approx(12617,
12617)$ 조건에 구식 공식을 대입하면 $-84.7\dB$가 나와 v3의 수치를 거의 그대로
재현한다. 교정된 공식으로 다시 스캔한 결과 이상 source-limit 전역 최적은
$\Delta=-2.10\GHz$, $T=130\degC$, $\xi=-31.3454\dB$로, v3보다 $50\dB$ 가까이
얕아졌고 온도 좌표도 $142.5\degC\to130\degC$로 이동했다.

반면 현실적 손실을 포함한 finite-loss 결과와 문헌 레퍼런스 재현은 구조적으로는
v3와 같은 결론을 유지한다. 다만 정량값은 갱신되었다: 실효 최적점은 여전히
$T\approx120$--$135\degC$ 부근이며, 앱 자체의 Ultra 프리셋으로 읽은 Sim 동작점
재현값은 $\xi\approx-7.15\dB$이다(문헌 $-7.8\dB$와 같은 자릿수).
\end{abstract}

\tableofcontents
\bigskip

%==================================================================
\section{v3 대비 달라진 것}
%==================================================================
이 절은 v3를 읽은 독자를 위한 변경 요약이다. 물리적 배경과 전체 논증은 이후
절에서 다시 설명한다.

\begin{enumerate}[nosep]
\item \textbf{(A) 이상 세기차 노이즈 공식 교정.} 무손실 트윈빔에서 광자수차
      $N_s-N_c$는 보존량이므로, 일반식은
      $S_{\mathrm{ideal}}=(G_s-G_c)^2/(G_s+G_c)$이고 $G_s-G_c=1$인 무손실
      한계에서 정확히 $1/(2G-1)=1/(G_s+G_c)$로 환원된다(Sim \textit{et al.}의
      레퍼런스 법칙과 일치). 구식 공식은 $G=2$에서만 이 식과 일치하고 고이득에서
      과도하게 깊은 squeezing을 주었다.
\item \textbf{(E) 밀도/온도 의존 충돌 결어긋남 도입.} 광학 결맞음에는 Rb
      self-broadening($\Gamma_{\mathrm{eff}}=\Gamma+\beta N$, AutoOD 검증된
      계수), 바닥상태 Raman 결맞음에는 transit-time floor + Rb--Rb 충돌항을
      추가했다. 그 결과 finite-loss squeezing이 $\sim121$--$131\degC$에서
      정점을 찍고 고온에서 감소하는, 문헌이 보고하는 거동을 재현한다.
\item \textbf{twin-beam gap 가드 신설.} 교정된 공식은 $G_s-G_c=1$이 성립하는
      진짜 파라메트릭 이득 영역에서만 엄밀하다. 이 관계가 깨지는(공명에서 먼
      배경 영역, 또는 고이득 근처의 수치적 불안정 영역) 지점에서는
      $(G_s-G_c)^2\to0$이 되어 실제 이득 없이도 가짜로 완벽한 squeezing이
      보고될 수 있다. 이번 스캔에서는 $G_s-G_c\in[0.5,1.5]$를 신뢰 구간으로 두고
      이를 벗어나는 $\delta$-argmin 후보를 제외한다.
\item \textbf{실행 방식.} 이번 이상 스캔(19389점)은 실행 환경의 장시간
      백그라운드 프로세스 제약으로 인해 T값 69개로 쪼개어 순차 실행 후 병합했다
      (부록 참조). 결과 데이터는 단일 실행이었다면 얻었을 것과 동일하다
      (결정론적 계산).
\end{enumerate}

%==================================================================
\section{질문의 분리 (v3과 동일)}
%==================================================================
이 보고서에서 서로 다른 세 질문을 분리한다.

\begin{enumerate}[nosep]
\item \textbf{실험 레퍼런스 재현}: Sim \textit{et al.}과 Lett 계열 문헌의 85Rb D1
      double-$\Lambda$ FWM squeezing 동작점을 GABES가 재현하는가?
\item \textbf{현실적 최효율 전선}: 검출효율, optical loss, in-cell 손실을 포함했을 때,
      최적에 거의 도달하면서 온도/이득/OD를 최소화하는 robust operating point는 어디인가?
\item \textbf{이상 source-limit}: \texttt{QE=1}, \texttt{loss=0}으로 검출 floor를 제거하면,
      Ultra 모형 자체가 허용하는 가장 깊은 세기차 스퀴징은 어디인가?
\end{enumerate}

%==================================================================
\section{좌표계와 레퍼런스 설정}
%==================================================================

\subsection{GABES detuning convention}
\begin{itemize}[nosep]
\item $\Delta$: 펌프 one-photon detuning, 기준은 $^{85}$Rb D1의 $F=2\to F'=3$ 선.
\item $\omega_{\mathrm{pump}}=\omega(F=2\to F'=3)+\Delta$.
\item 표준 seeded FWM readout은 $(-)$ Raman branch.
\item 바닥 초미세 분리 $\nu_{\mathrm{HF}}=3.035732439\GHz$.
\end{itemize}

Sim \textit{et al.}의 실험 최적점 $\Delta\approx+0.9\GHz$도 GABES와 동일하게
$F=2\to F'=3$ 기준이다(Sim Fig.~2 캡션: ``set at 0.9~GHz towards the
blue-detuning from the $^{85}$Rb $5S_{1/2}$, F=2 to $5P_{1/2}$, F=3
transition''). 따라서 이번 스캔의 이상 최적 $\Delta=-2.10\GHz$(v3과 동일 좌표)와
Sim의 $+0.9\GHz$는 기준점 차이로 설명되는 것이 아니라, 같은 기준에서 실제로
서로 다른 detuning을 가리키는 별개의 지점이다. 이 차이의 이유는 목적함수가
다르기 때문이다(무손실 ideal source-limit 최적 vs 실측 finite-loss IDS
최적) — 자세한 논의는 문헌 비교 절의 ``일치하지 않는 부분과 이유'' 참조.

\subsection{고정 파라미터}
\begin{table}[h]
\centering
\caption{seeded 85Rb D1 FWM 레퍼런스 설정 (v3과 동일).}
\label{tab:refparams}
\begin{tabular}{lll}
\toprule
항목 & 값 & 비고\\
\midrule
$P_{\mathrm{pump}}$ & $600\,\mathrm{mW}$ & Sim 레퍼런스\\
$P_{\mathrm{seed}}$ & $8\,\mu\mathrm{W}$ & Sim 레퍼런스\\
line-strength residual & $0.74$ & A/E 교정 후에도 재보정 불필요\\
셀 길이 $L$ & $12.5\,\mathrm{mm}$ & GABES 상수\\
빔 허리 $w_{\mathrm{pump}}/w_{\mathrm{seed}}$ & $530/330\,\mu\mathrm{m}$ & GABES 상수\\
branch & $-1$ & standard seeded FWM branch\\
\bottomrule
\end{tabular}
\end{table}

현실 조건 스캔은 $\mathrm{QE}=0.92$, post-cell optical loss $5.5\%$를 사용해
$\eta=0.8694$이다. 이상 source-limit 스캔은 $\mathrm{QE}=1$, post-cell loss
$0$으로 둔다. 이때 detection floor는 물리적으로 $-\infty$이며, 코드의
guardrail 표시는 $-3000\dB$이다.

%==================================================================
\section{현실 손실 포함 Ultra 결과: 레퍼런스와 맞는가}
%==================================================================

\subsection{finite-loss Ultra frontier}
$\Delta\in[-3.5,+3.0]\GHz$, $T\in[60,150]\degC$ (v2/v3과 동일 그리드)에서
교정된 물리로 재실행한 결과는 그림~\ref{fig:frontier_v4}와 같다.

\begin{figure}[p]
\centering
\includegraphics[width=\textwidth]{squeezing_frontier_v4.png}
\caption{finite-loss Ultra 스캔(v4, 교정된 물리). 왼쪽: $\xi(\Delta,T)$ heatmap. 가운데:
최적에 $0.1\dB$ 이내로 도달하는 최저 온도 전선. 오른쪽:
최효율 detuning에서의 $\xi(T)$.}
\label{fig:frontier_v4}
\end{figure}

\begin{table}[h]
\centering
\caption{finite-loss Ultra 스캔의 핵심 좌표 (v4, 교정된 물리).}
\label{tab:v4summary}
\begin{tabular}{lrrrrr}
\toprule
점 & $\Delta$ [GHz] & $T$ [$^\circ$C] & $\xi$ [dB] & $G_s$ & in-cell OD\\
\midrule
전역 최적 & $-2.20$ & $135$ & $-8.8342$ & $635.6$ & $0.0000$\\
최효율점($0.1\dB$ 이내) & $-2.10$ & $125$ & $-8.7665$ & $53.7$ & $0.0002$\\
Sim 동작점 재현(앱 Ultra 프리셋) & $+0.90$ & $121$ & $-7.1493$ & $14.1$ & $0.0439$\\
\bottomrule
\end{tabular}
\end{table}

detection floor $10\log_{10}(1-\eta)=-8.8406\dB$에 전역 최적이 사실상 도달한다
($+0.0063\dB$ 차이). 최효율점은 $T\approx125\degC$의 훨씬 작은 이득
영역($G_s\approx54$)에서 최적보다 $0.068\dB$ 얕을 뿐이다.

\textbf{Sim 동작점 재현값이 v3과 달라진 이유.} v3은 $\xi\approx-7.88\dB$를
보고했으나, A(공식)+E(충돌 결어긋남) 교정 후 같은 조건(앱의
\texttt{FWMScheme.compute()} 경로, $\Delta=0.9\GHz$, $T=121\degC$,
$\delta=-8\MHz$, Ultra 프리셋)에서 $\xi=-7.1493\dB$로 재현된다. 이는 여전히
문헌 $-7.8\dB$와 같은 자릿수이며, 온도 스케일($T\sim120\degC$)과 이득
스케일($G_s\approx14$)도 여전히 일치한다. 차이는 이상 검출 floor를 향해
쏠리던 A 이전 공식의 왜곡이 제거되고, E가 이 부근에서 약한 추가 흡수를 더했기
때문으로, 재보정(line-strength residual 변경)은 필요하지 않았다.

\subsection{in-cell 손실과 충돌 결어긋남이 함께 만드는 내부 최적}
그림~\ref{fig:od_v4}는 in-cell optical depth를 보여준다. 공명 근처에서는
온도가 올라갈수록 gain과 OD가 함께 증가한다. E 도입 이후에는 추가로, 충돌
결어긋남이 매우 고온에서 gain 자체를 깎기 시작해 $\xi(T)$가 단순히 손실로만
제한되는 것이 아니라 이중의 고온 페널티(손실 + 결어긋남)를 받는다.

\begin{figure}[p]
\centering
\includegraphics[width=\textwidth]{squeezing_frontier_od_v4.png}
\caption{v4의 in-cell OD 진단(교정된 물리). 왼쪽: $\mathrm{OD}(\Delta,T)$ map.
오른쪽: 대표 detuning들의 $\xi(T)$ slice.}
\label{fig:od_v4}
\end{figure}

%==================================================================
\section{이상 검출 한계 스캔: \texttt{QE=1}, \texttt{loss=0}}
%==================================================================

\subsection{실행 조건}
v3과 동일한 넓고 촘촘한 격자를 교정된 물리로 재실행했다.

\begin{table}[h]
\centering
\caption{v4 이상 source-limit 스캔 설정 (v3과 동일 그리드).}
\label{tab:idealrun}
\begin{tabular}{ll}
\toprule
항목 & 값\\
\midrule
fidelity & Ultra\\
$\mathrm{QE}$, post-cell loss & $1.0$, $0$\\
$\Delta$ 범위 & $[-6,+8]\GHz$, step $0.05\GHz$\\
$T$ 범위 & $[50,220]\degC$, step $2.5\degC$\\
$\delta$ scan window & $\pm1.0\GHz$\\
coarse points & $201$\\
velocity grid & $2.0\,\mathrm{m/s}$, cutoff $4\sigma$\\
총 격자 & $281\times69=19389$ points\\
실행 방식 & T값 69개로 분할, 개별 순차 실행 후 병합 (부록)\\
\bottomrule
\end{tabular}
\end{table}

고온/고이득 일부 외곽점에서 segment propagation의 overflow warning이
발생했지만(코드 내 clip으로 처리), 최종 병합 배열의 $\xi$, $G_s$, $G_c$에는
NaN/Inf가 0개였다. $\delta$ window edge에 걸린 점은 $107/19389$개, twin-beam
gap 가드에 걸린 점은 $0/19389$개였다. 전역 최적은 두 제외 조건 모두에
해당하지 않는 내부점이다.

\subsection{전역 최적 결과 (교정된 물리)}
\begin{figure}[p]
\centering
\includegraphics[width=\textwidth]{squeezing_frontier_ideal_qe1_v4.png}
\caption{v4 이상 검출 한계 스캔(교정된 물리). 왼쪽: $\xi(\Delta,T)$ heatmap과
전역 최적(청록 원). 가운데: 각 $T$에서 얻을 수 있는 최저 $\xi$. 오른쪽: 최적
detuning $\Delta=-2.10\GHz$에서의 $\xi(T)$ slice.}
\label{fig:ideal_qe1_v4}
\end{figure}

\begin{table}[h]
\centering
\caption{v4 이상 source-limit 최적점 (교정된 물리) vs v3 (구식 공식).}
\label{tab:idealopt}
\begin{tabular}{lrr}
\toprule
항목 & v4 (교정됨) & v3 (구식 공식)\\
\midrule
$\Delta$ & $-2.10\GHz$ & $-2.10\GHz$\\
$T$ & $130.0\degC$ & $142.5\degC$\\
$\delta$ & $-320\MHz$ & $-120\MHz$\\
$\xi$ & $-31.3454\dB$ & $-81.9469\dB$\\
$G_s$ & $415.3$ & $12617.07$\\
$G_c$ & $414.8$ & $12615.04$\\
$G_s-G_c$ (gap) & $+0.540$ & $\approx0$ (신뢰 불가)\\
in-cell OD & $0.0004$ & $0.000$\\
edge/gap-invalid? & 아니오 & (가드 없음)\\
\bottomrule
\end{tabular}
\end{table}

%==================================================================
\section{왜 v3의 $-81.9\dB$는 물리가 아니라 공식 버그였는가}
%==================================================================

\subsection{같은 데이터, 두 공식}
v3이 보고한 좌표 근방($\Delta=-2.10\GHz$)의 v4 스캔 데이터에서 $T=142.5\degC$
행을 직접 확인하면 $G_s=6608$, $G_c=6607$인 지점이 있다(v3의 $G_s=12617$과
같은 자릿수). 이 \emph{동일한} $(G_s,G_c)$ 쌍에 두 공식을 적용하면:

\begin{equation}
S_{\mathrm{ideal}}^{\text{(구식)}}=\frac{G_s+G_c-2\sqrt{G_sG_c-1}}{G_s+G_c}
\ \Longrightarrow\ 10\log_{10}S^{\text{(구식)}}=-78.94\dB
\end{equation}
\begin{equation}
S_{\mathrm{ideal}}^{\text{(교정)}}=\frac{(G_s-G_c)^2}{G_s+G_c}
\ \Longrightarrow\ 10\log_{10}S^{\text{(교정)}}=-29.51\dB
\end{equation}

구식 공식의 결과($-78.94\dB$)는 v3의 $-81.9469\dB$와 거의 일치한다(잔여
차이는 v3과 v4의 정확한 $\delta$/$T$/coarse 해상도가 미세하게 다르기
때문이다). 전체 병합 격자에서 $G_s$가 가장 큰 점($\Delta=-2.05\GHz$,
$T=150\degC$, $G_s\approx12830$, v3의 $12617.07$과 사실상 동일)에 구식
공식을 대입하면 $-84.72\dB$가 나와 v3의 수치를 더욱 정밀하게 재현한다.

\subsection{왜 구식 공식이 고이득에서 터지는가}
구식 공식 $(G_s+G_c-2\sqrt{G_sG_c-1})/(G_s+G_c)$는 $G_s\approx G_c\gg1$인
극한에서 $\sqrt{G_sG_c-1}\to\sqrt{G_sG_c}$로 근사되고, 분자가 두 개의 거의
같은 큰 수의 차(catastrophic cancellation)가 되어 부동소수점 오차와
결합하면서 물리적으로 무의미한 초저 노이즈 값을 만든다. 반면 교정식
$(G_s-G_c)^2/(G_s+G_c)$는 애초에 $G_s-G_c$라는 \emph{작은} 차이를 직접
제곱하므로 이런 수치적 붕괴가 없고, $G_s=G_c$일 때 정확히 $0$을 준다 — 이는
무손실 극한에서 $G_s-G_c=1$이 성립해야 한다는 물리적 제약과 어긋나는 신호로,
바로 이 지점(공명에서 먼 배경 영역 또는 수치적으로 불안정한 고이득 영역)이
twin-beam gap 가드가 걸러야 하는 영역이다.

\subsection{v4 전역 최적이 물리적으로 신뢰할 수 있는 이유}
v4의 전역 최적($\Delta=-2.10\GHz$, $T=130\degC$)은 $G_s-G_c=+0.540$으로
gap 가드 구간 $[0.5,1.5]$ 안에 있고, 인접 온도($T=127.5\degC$: $\xi=-28.40\dB$,
$T=132.5\degC$: $\xi=-26.76\dB$)에서도 매끄럽게 변하는 내부 극값이다. 이는
검증된 Sim 레퍼런스 지점의 $\delta$ 전체 스캔에서 gap이 $[0.55,0.98]$
범위를 유지하는 것과 같은 성격의 신뢰 가능한 파라메트릭 영역이다.

%==================================================================
\section{문헌 레퍼런스와의 비교}
%==================================================================

\subsection{85Rb D1 double-$\Lambda$ squeezing 최적점}
\begin{table}[h]
\centering
\caption{문헌과 GABES v4 결과의 비교.}
\label{tab:literature}
\begin{tabular}{p{3.3cm}p{3.4cm}p{3.0cm}p{4.8cm}}
\toprule
출처 & 보고된 조건 & squeezing/gain & GABES v4와의 관계\\
\midrule
McCormick \textit{et al.} (2007) &
hot Rb vapor FWM &
$-3.5\dB$ measured, $-8.1\dB$ loss-corrected &
GABES finite-loss 결과의 크기($-8\dB$대)와 같은 scale.\\
\addlinespace
Pooser \textit{et al.} (2009) &
85Rb D1, $\Delta\approx+0.8\GHz$ cut, 400 mW pump &
$-8.0(2)\dB$; detection efficiency $\approx0.90$라 $-10\dB$ 상한 &
finite-loss GABES가 $\eta$ floor에 의해 $-8$에서 $-10\dB$대로 제한되는 것과 일치.\\
\addlinespace
Sim \textit{et al.} (2025) &
85Rb D1, $T=121\degC$, $\Delta\approx+0.9\GHz$, $\delta=-8\MHz$ &
gain $\approx15$, IDS $-7.8\dB$, slope estimate $-8.7\dB$ &
GABES v4 Sim 동작점 재현 $\xi\approx-7.15\dB$, $G_s\approx14$로 여전히 같은
자릿수(v3의 $-7.88\dB$보다는 다소 벌어짐, 4절 참조).\\
\addlinespace
Allen \textit{et al.} (2026) &
85Rb 12.5 mm cell, $T\sim118\degC$ near optimized detuning &
long-term stable IDS $-7.8\dB$ &
$T\sim118$--$135\degC$ 부근이 robust operating point임을 뒷받침(v4 최효율점
$T=125\degC$, 전역 최적 $T=135\degC$).\\
\addlinespace
Jasperse \textit{et al.} (2011) &
85Rb FWM with distributed gain/loss model &
gain과 absorption 경쟁을 analytic model로 설명 &
GABES Ultra의 in-cell loss + 충돌 결어긋남(E) 채널과 같은 물리 방향.\\
\bottomrule
\end{tabular}
\end{table}

\subsection{일치하는 부분}
\begin{enumerate}[nosep]
\item \textbf{온도 scale}: Sim/Allen의 $118$--$121\degC$와 GABES 최효율점
      $125\degC$가 일치한다.
\item \textbf{squeezing scale}: Sim의 $-7.8\dB$와 GABES v4 재현값
      $-7.15\dB$가 같은 자릿수다.
\item \textbf{gain scale}: Sim의 gain $\approx15$와 GABES v4의
      $G_s\approx14$가 같은 order이다.
\item \textbf{손실 제한}: detection/intrinsic loss 제한이 GABES의 $\eta$
      floor 및 in-cell OD penalty와 같은 방향이다.
\item \textbf{고온 한계}: E 도입 이후, 온도를 계속 올리면 collision
      dephasing이 squeezing을 악화시킨다는 문헌의 정성적 주장을 GABES가 처음으로
      재현한다(finite-loss $\xi(T)$가 $\sim121$--$135\degC$에서 정점).
\end{enumerate}

\subsection{일치하지 않는 부분과 이유}
v4의 ideal source-limit 최적점은 여전히 문헌의 실험 최적점과 일치하지
\emph{않는다} — 이는 v3에서와 같은 이유(목적함수 차이: 손실 없는
source-limit vs 실측 IDS)이며, 교정 후에도 그 원칙은 바뀌지 않는다. 다만
수치는 훨씬 얌전해졌다: $-81.9\dB\to-31.3\dB$, $G_s\,12617\to415$. 여전히
$G_s\sim400$은 실험적으로 검증되지 않은 초고이득 영역이므로, 이 좌표를 정량
실험 예측으로 쓰려면 그 영역의 별도 calibration과 완전한 noise model이
필요하다는 v3의 결론은 유지된다.

%==================================================================
\section{물리적 해석}
%==================================================================

\subsection{finite-loss frontier가 여전히 중요한 이유}
$S(\eta)=\eta S_{\mathrm{cell}}+(1-\eta)$이므로 $\eta=0.8694$에서 floor는
$-8.8406\dB$이다. $G_s\gtrsim54$면 이미 floor에 거의 도달하므로, 더 큰 이득은
비용(pump scatter, 충돌 결어긋남 등 excess noise)만 키운다. 실험 권고점은
"가장 깊은 점"이 아니라 "floor에 거의 도달하는 가장 낮은 온도와 이득의
점"이다 — 이는 A/E 교정과 무관하게 v3부터 유지되는 원칙이다.

\subsection{교정 후 이상 스캔이 여전히 $-30\dB$대인 이유}
무손실 balanced twin beam에서 $G_c\simeq G_s-1$이면
$S_{\mathrm{ideal}}=1/(2G_s-1)$이다. v4 최적의 $G_s\simeq415$를 대입하면
\begin{equation}
10\log_{10}\!\left(\frac{1}{2\times415-1}\right)\approx-29.2\dB
\end{equation}
로, 실제 값 $-31.3454\dB$와 같은 order다(차이는 이득 불균형과 propagation
detail을 포함하기 때문). 이는 v3이 인용했던 $5/(8G^2)$ 스케일링(구식 공식의
점근형)과는 다른, $1/(2G)$ 스케일링이며 — \emph{이 지수의 차이 자체가 A
교정의 핵심}이다.

\subsection{실험 안티스퀴징은 손실만으로 설명되지 않는다}
손실은 squeezed difference noise를 shot noise 쪽으로 되돌릴 뿐 $0\dB$ 위로
올리지 못한다. 실험에서 difference channel이 양의 dB로 올라가거나 predicted
loss floor보다 훨씬 얕다면 원인은 pure loss가 아니라 excess noise다. E
교정으로 이제 충돌 결어긋남이 고온에서 실제로 gain을 깎는 채널로 모델에
들어와 있지만, 이는 여전히 pump scatter나 anti-squeezed sum noise leakage
같은 완전한 noise model의 일부만 포함한다.

%==================================================================
\section{결론}
%==================================================================

\begin{enumerate}[nosep]
\item \textbf{v3의 이상 스캔 결과는 공식 버그였다}: $-81.9469\dB$는 물리적
      실체가 아니라 이상 노이즈 공식의 catastrophic-cancellation 버그였다.
      같은 $(G_s,G_c)$에 구식/신식 공식을 각각 적용하면 $-84.7\dB$(구식,
      v3과 일치) vs $-21.4\dB$(신식)로 확인된다.
\item \textbf{교정된 이상 source-limit}: $\Delta=-2.10\GHz$, $T=130\degC$,
      $\xi=-31.3454\dB$, $G_s\approx415$ — 여전히 실험 조건과는 다른
      영역이지만, 훨씬 덜 극단적이고 twin-beam gap 가드로 신뢰성이
      확인되었다.
\item \textbf{finite-loss 결과는 구조적으로 유지되며 정량적으로 갱신됨}:
      최효율점 $\Delta\approx-2.1\GHz$, $T\approx125\degC$; Sim 동작점
      재현은 $\xi\approx-7.15\dB$로 문헌과 같은 자릿수를 유지한다.
\item \textbf{충돌 결어긋남(E) 도입으로 고온 한계가 처음 재현됨}: squeezing이
      온도에 따라 무한정 좋아지지 않고 $\sim121$--$135\degC$에서 정점을
      찍은 뒤 악화되는, 문헌이 강조해온 물리가 모델에 나타난다.
\item \textbf{일반 원칙}: 물리 공식을 교정할 때는 그 공식을 소비하는 모든
      다운스트림(스킬 스크립트, 분석 스크립트, 이전 보고서의 수치)을
      재검증해야 한다. 부정확한 공식이 우연히 가려주던 정의역-이탈 지점을
      교정된 공식이 적나라하게 드러낼 수 있다 — 이번 twin-beam gap 가드가
      바로 그 사례다.
\end{enumerate}

\noindent\textbf{한 줄 요약.} \emph{v3의 $-81.9\dB$ "이상 한계"는 공식
버그였다. 교정 후 GABES는 문헌의 85Rb D1 FWM squeezing 동작점을 현실 손실
조건에서 여전히 잘 재현하며($\xi\approx-7.15\dB$), 이상 source-limit은 훨씬
덜 극단적인 $-31.3\dB$/$G_s\approx415$로 재확정되었다.}

\appendix
%==================================================================
\section{수치 재현 정보}
%==================================================================

\begin{itemize}[nosep]
\item finite-loss v4 산출물:
      \texttt{analysis/squeezing/squeezing\_frontier\_finite\_loss\_v4/squeezing\_frontier.npz},
      \texttt{docs/squeezing\_report/squeezing\_frontier\_v4.png},
      \texttt{docs/squeezing\_report/squeezing\_frontier\_od\_v4.png}
      (\texttt{docs/squeezing\_report/make\_v4\_finite\_loss\_figs.py}로 생성).
\item ideal v4 산출물:
      \texttt{analysis/squeezing/squeezing\_frontier\_ideal\_v4\_chunks/T\_*/squeezing\_frontier.npz}
      (T값 69개, 각 전체 $\Delta$축$\times$1개 T), 병합 결과
      \texttt{analysis/squeezing/squeezing\_frontier\_ideal\_v4\_merged.npz}
      (\texttt{docs/squeezing\_report/merge\_v4\_ideal\_chunks.py}로 생성).
\item ideal v4 그림:
      \texttt{docs/squeezing\_report/squeezing\_frontier\_ideal\_qe1\_v4.png}.
\item ideal v4 finite check: 병합 격자 $19389$점 중 $\xi$, $G_s$, $G_c$
      NaN/Inf $0$개. $\delta$-edge 제외 $107$개, twin-beam gap 가드 제외
      $0$개. 전역 최적은 두 제외 조건 모두 해당하지 않는 내부점.
\item \textbf{실행 방식(장시간 백그라운드 제약)}: 실행 환경에서 단일
      장시간(수 시간) 백그라운드 프로세스가 하네스/OS 레벨 추적과 무관하게
      중도 종료되는 현상이 관측되었다(원인 불명, 코드 문제 아님). 이를
      우회하기 위해 전체 $281\times69$ 격자를 T값 69개로 쪼개어(각 청크 =
      전체 $\Delta$축 $\times$ 1개 T, $\sim$180--210초) 순차 실행하고 완료된
      npz를 병합했다. 계산은 결정론적이므로 단일 실행 결과와 동일하다.
\item \textbf{공식 검증}: \texttt{gabes/observables.py}의
      \texttt{ideal\_twin\_beam\_noise(G\_s, G\_c)}가 교정된 식을 구현한다.
      \texttt{.claude/skills/fwm-squeezing-frontier/scripts/scan\_squeezing\_frontier.py}의
      \texttt{squeezing\_map()}이 이를 직접 호출하도록 수정되었다(이전에는
      구식 공식을 자체 재구현하고 있어 A 교정이 반영되지 않았었다).
\end{itemize}

%==================================================================
\section{참고문헌}
%==================================================================

\begin{enumerate}[label={[\arabic*]},leftmargin=*]
\item G. Sim, H. Kim, H. S. Moon, ``Intensity-difference squeezing from four-wave mixing
      in hot $^{85}$Rb and $^{87}$Rb atoms in single diode laser pumping system,''
      \textit{Scientific Reports} \textbf{15}, 7727 (2025).
      \url{https://www.nature.com/articles/s41598-025-86479-w}
\item C. F. McCormick, V. Boyer, E. Arimondo, P. D. Lett,
      ``Strong relative intensity squeezing by four-wave mixing in rubidium vapor,''
      \textit{Optics Letters} \textbf{32}, 178--180 (2007).
      \url{https://doi.org/10.1364/OL.32.000178}
\item R. C. Pooser, A. M. Marino, V. Boyer, K. M. Jones, P. D. Lett,
      ``Quantum correlated light beams from non-degenerate four-wave mixing in an atomic vapor:
      the D1 and D2 lines of $^{85}$Rb and $^{87}$Rb,''
      \textit{Optics Express} \textbf{17}, 16722--16730 (2009).
      \url{https://doi.org/10.1364/OE.17.016722}
\item M. Jasperse, L. D. Turner, R. E. Scholten,
      ``Relative intensity squeezing by four-wave mixing with loss: an analytic model and
      experimental diagnostic,'' \textit{Optics Express} \textbf{19}, 3765--3774 (2011).
      \url{https://arxiv.org/abs/1012.3482}
\item C. H. Allen, M. San Soucie, R. J. Benson, O.-M. Glezakou-Elbert, R. Jimenez,
      J. Morrell-Falvey, Y.-Z. Ma,
      ``Long-term stabilization of intensity-difference squeezing from four-wave mixing in
      rubidium vapor,'' \textit{Optics Express} \textbf{34}, 19382--19399 (2026).
      \url{https://doi.org/10.1364/OE.600911}
\end{enumerate}

\end{document}
